% !TeX spellcheck = hu_HU
% !TeX encoding = UTF-8
% !TeX program = xelatex
%----------------------------------------------------------------------------
\chapter{Továbbfejlesztési lehetőségek}
%----------------------------------------------------------------------------

\section{Javítások ajánlása}

A jelenlegi eszköz elsődlegesen a különböző szabályoktól való eltérést hivatott segíteni megkeresni.
Ennek ellenére, a Clang támogat úgynevezett ``FixIt'' segítségeket.
Ezek a normál diagnosztikával egyszerre tudnak megjelenni és konkrét javítási lehetőségeket tudnak javasolni a programozó számára.

Ezt a lehetőséget felhasználva, a blAST eszköz képes lenne nem csak a hibákat megjeleníteni, hanem egyes esetekben a konkrét megoldást is képes lehet kiírni a felhasználó számára.
Ezzel lehetővé válik sokkal egyértelműbb hiba üzeneteket létrehozni, amik javítják a felhasználási élményt.

Ezt vélhetőleg a legjobban egy kibővített C4 ``diagnose'' függvénnyel lehetséges megvalósítani, amelyik a ``fixit'' paramétert is értelmezni képes.
Például \aref{lst:blast-fixit}.~lista egy lehetséges szintaktikai megoldást mutat be.
A megtalált kifejezés helyére a megadott szöveget (``something'') fogja felajánlani, mint javított kódrészlet.
Valós példákban vélhetőleg a kapott találat szöveges reprezentációjának valamilyen módosításaként lenne érdemes a javítást összeállítani.

\begin{lstlisting}[label=lst:blast-fixit,caption={FixIt segítségek blAST-al}]
blast ast
	match into callback { |diag|
		diagnose diag fixit "something"
	}
	such call expression
	that # ...
\end{lstlisting}

\section{Több nyelv támogatása}

A jelenlegi rendszer specifikusan a Clang alapú C és C++ nyelveket célozza meg.
Emiatt nagy része a megvalósított C4 szintű elemeknek ezen nyelvek sajátosságait célozzák meg, és sok helyen feltételeznek olyan dolgokat, amelyek specifikusan csak ebben a felállásban igazak.

Ennek ellenére, a legkülsőbb rétege a blAST eszköznek nem követeli meg, hogy ezeket a nyelveket használja.
A konkrét {\tt such ... that ... } szintaxis például teljesen független minden másik elemtől, mert csak a paramétereinek megfelelő struktúráját hivatott biztosítani.
Az ASTMatcher kifejezés előállítása után sem tartozik hozzá nyelv specifikus elem: az összerakott kifejezés egy egyszerű sztring formájában utazik a rendszerben, amíg a Clang könyvtárak azt fel nem dolgozzák.

A jövőre való tekintettel, érdemes lehet megnézni, mennyire lehet a többi részét más nyelvekhez illeszteni.
Például a Java és C\# nyelvek érdekesek lehetnek nagy elterjedtségük miatt.
Utóbbi esetén a Roslyn \cite{microsoftDotnetRoslyn2026} fordító a Clang könyvtárakhoz hasonló módon sok kapcsolódási pontot biztosít egy C\# projekt kódjának értelmezéséhez, kezeléséhez.
Ez könnyítheti annak beépítését a blAST keretbe.

Nehézséget a AST-ben kereső rendszer megvalósítása okozhat. 
Az ASTMatchers könyvtár szöveges minták támogatása miatt nagyon könnyen használható, ha hasonló támogatás nem létezik az említett nyelvekhez, akkor a kereső kifejezések előállítása jelentősen bonyolultabb lesz.
Ez azonban nem megoldhatatlan, csak a megfelelő elemeket direkt módon összerakó függvényeket kell a C4 nyelv számára elérhetővé tenni, és akkor a többi elem összerakását a jelenlegi blAST/Clang módszerrel ekvivalensen meg lehet csinálni.

